% Cote 119: offprints, reports, courses, CV and letters (typed and printed part): archive cote 119, site page 265 (PDF page 266).
% typescript. Draft compared: none.
% AI-transcribed by Claude Opus 5.5 (claude-opus-5-5, adaptive thinking, effort high), Message Batches API, batch r119_full (msgbatch_01DwZNm5Y7YsKBWyAN4sFs8w), 2026-11-13; not yet checked by eye.
\scanpage{119/265}
\mspage{12}

sous la pression du manque d'enseignants. Il faudrait que le docteur mention
``recherche'' ait une personnalité scientifique assez forte pour être jugé capable
de former lui-même des chercheurs de valeur. Quant au doctorat mention ensei-%
gnement, il serait décerné sur critères pédagogiques, se basant sur un minimum
de quatre à cinq années de pratique de l'enseignement comme assistant et maître
assistant, avec pendant une année au moins la responsabilité complète d'un
cours. De plus, les activités d'enseignement du candidat devront avoir porté
sur au moins trois sujets différents. Il serait tenu compte, dans l'attri-%
bution de ce doctorat, des avis des chefs de service sous la direction desquels
le candidat aura été assistant, ainsi que des avis exprimés année par année par
les étudiants enseignés par le candidat. Un étudiant, élu parmi ceux ayant été
enseignés par le candidat, pourrait utilement faire partie du jury décidant de
l'attribution du titre de docteur d'enseignement.

4. Pour que la revalorisation de la vocation d'enseignant à l'Université soit
authentique, et ne se réduise à une façon pieuse de cacher une dévalorisation
de la fonction de maître de conférences ou de professeur à l'Université, qu'il
faut au contraire éviter à tout prix, il est essentiel que le doctorat de
Mathématique (mention enseignement) ne soit décerné que sur des critères de
sélection très sérieux, et qu'il ne puisse être question pour son attribution
d'invoquer des arguments d'\uline{ancienneté} du candidat dans les fonctions d'assis-%
tant et de maître assistant. Ne devrait être attribué ce titre qu'à un candidat
ayant les qualifications voulues pour devenir un bon maître de conférences
(mention enseignement), i.e. essentiellement, un candidat ayant prouvé ses capa-%
cités pour assumer la responsabilité complète d'un cours, et donnant des
garanties d'ouverture intellectuelle suffisantes pour qu'on soit fondé à penser
qu'il sera en mesure d'assumer la responsabilité complète de n'importe quel
cours de Mathématique d'un niveau pré-licence. Pour fonder cette conviction,
en plus des appréciations des étudiants et des chefs de service sur les acti-%
vités d'enseignement proprement dites, le jury pourra tenir compte de tous
autres signes d'initiative intellectuelle du candidat, que ce soit dans la
thèse de celui-ci, ou par sa participation active à des séminaires, la rédac-%
tion de notes polycopiées, la publication d'articles de recherche mathématiques

\margin{265}
